% Generated draft: requires mathematical and source review.
\documentclass[11pt]{article}
\usepackage[margin=25mm]{geometry}
\usepackage{fontspec}
\setmainfont{Latin Modern Roman}
\usepackage{amsmath,amssymb,mathtools,unicode-math}
\setmathfont{Latin Modern Math}
\usepackage{enumitem,needspace,xurl,hyperref}
\setlength{\parindent}{0pt}
\setlength{\parskip}{5pt}
\setlength{\emergencystretch}{3em}
\newcommand{\N}{\mathbb{N}}
\newcommand{\Z}{\mathbb{Z}}
\newcommand{\Q}{\mathbb{Q}}
\newcommand{\R}{\mathbb{R}}
\newcommand{\C}{\mathbb{C}}
\newcommand{\F}{\mathbb{F}}
\newcommand{\A}{\mathbb{A}}
\newcommand{\PP}{\mathbb{P}}
\newcommand{\E}{\mathbb{E}}
\newcommand{\eps}{\varepsilon}
\newcommand{\dd}{\,\mathrm{d}}
\begin{document}
\section*{UnsolvedMath ID: 30000385}
\textbf{Source problem code:} catalogue entry
\subsection{Definitions and mathematical statement}
Let \(G\) be a connected graph. It is locally finite if every vertex has finite degree, planar if it has a crossing-free planar embedding, and \(4\)-connected if it has at least five vertices and deleting any set of at most three vertices leaves it connected. Use the conventional finite-simple-graph terminology extended to locally finite graphs, without a restriction on the number of ends.

For a locally finite connected graph, let \(|G|\) denote its Freudenthal compactification: regard each edge as a topological interval and adjoin the graph's ends. An end is an equivalence class of rays, where two rays are equivalent when no finite vertex set separates their tails. Neighborhoods of an end are given by components remaining after finite vertex deletion, together with their corresponding ends and appropriate incident edge intervals.

A Hamilton circle is a subspace \(C\subseteq|G|\) homeomorphic to \(S^1\) and containing every vertex of \(G\). Equivalently there is a topological embedding \(\gamma:S^1\hookrightarrow|G|\) with \(V(G)\subseteq\gamma(S^1)\). The source conjecture is
\[
\forall G\quad
\bigl(G\text{ planar, locally finite, and }4\text{-connected}\bigr)
\ \Longrightarrow\ 
\exists C\subseteq|G|\text{ a Hamilton circle}.
\]
For finite graphs this specializes to a Hamilton cycle; for infinite graphs a finite cycle or a spanning double ray is not a substitute for the required topological circle. No one-endedness, finite-end assumption, or stronger planar embedding condition is added.

The supplied literature assessment checked 2026-08-21 says that the general conjecture remains open, with substantial published partial cases but a remaining hypothesis gap. It does not reproduce the exact partial-case conditions.
\subsection{Short English statement}
Does every locally finite, \(4\)-connected planar graph contain a topological circle through all vertices in its end compactification? The conjecture includes arbitrary end structure, not only one-ended graphs.
\subsection{Sources}
\begin{itemize}
\item Local UnsolvedMath registry, permanent ID 30000385, catalogue entry, \emph{Hamilton Circles in Infinite Planar Graphs}.
\item \emph{Combinatorics} (2006), Oberwolfach Reports, pp. 5--72, DOI \url{https://doi.org/10.4171/owr/2006/01}.
\item Supplied reference labeled \emph{Hamilton cycles in planar locally finite graphs}: \url{https://doi.org/10.1137/050631458}.
\item Supplied statement-verification claim and literature triage checked 2026-08-21; no full author list or partial-result hypotheses are supplied.
\end{itemize}
\end{document}
